\chapter{Income and Substitution Effects}\label{ch:b1:income-and-substitution-effects}

In 1991 Sweden began to tax the carbon in fuels, at 250 kronor a tonne of
carbon dioxide, about 30 dollars. % ledger: b1:sweden-co2-tax
The tax was part of a wider reform that cut the top marginal rate of income
tax from $80\%$ to $50\%$. % ledger: b1:sweden-1991-reform
Fuel became dearer and, for many households, take-home pay rose. If the
money comes back, why would anyone drive less? The answer splits every
price change in two: a change in what the goods cost relative to each
other, and a change in what the consumer's income can buy. The first
survives any cheque in the post.

\begin{openingchart}{Carbon dioxide from road transport per person, Sweden
and the unweighted mean of ten western European countries (Austria,
Belgium, Denmark, Finland, France, Germany, the Netherlands, Norway,
Switzerland, the United Kingdom), 1960--2022, tonnes per person. The
vertical line marks 1991.}{Community Emissions Data System (CEDS),
emissions by country and sector, road transport; World Bank, World
Development Indicators, population}
\begin{tikzpicture}
\begin{axis}[omtime, xlabel={Year}, ylabel={Tonnes of CO$_2$ per person},
  xmin=1960, xmax=2022, ymin=0, ymax=2.6, width=0.72\linewidth,
  height=0.48\linewidth, legend pos=south east]
  \addplot+[omDef, very thick, no marks] table[x=year, y=sweden] {figdata/bachelor-1/06-income-and-substitution-effects/out/road_co2.csv};
  \addplot+[omThm, thick, dashed, no marks] table[x=year, y=peers] {figdata/bachelor-1/06-income-and-substitution-effects/out/road_co2.csv};
  \draw[guide] (axis cs:1991,0) -- (axis cs:1991,2.6);
  \legend{Sweden, Mean of ten countries}
\end{axis}
\end{tikzpicture}
\end{openingchart}

Before 1991 Sweden's road emissions per person were higher than its
neighbours', though theirs grew faster (from 1960 to 1990 Sweden's rose
from 0.84 to 2.15 tonnes, the neighbours' mean from 0.51 to 1.93); % figdata: road_co2
from 1990 to 2005 Sweden's stayed near 2.2 tonnes while the
neighbours' mean rose from 1.93 to 2.23; by 2019 Sweden was at 1.55 tonnes
against 1.93. % figdata: road_co2_points
The chart alone cannot say how much of this the tax did. Fuel prices,
cars, biofuels and incomes all changed over sixty years; the record box
reports a study that builds a comparison for Sweden out of other
countries.

\begin{recall}
The consumer of \cref{ch:b1:consumer-choice} chooses the best bundle in her
\omterm{def:b1:consumer-choice:budget}{budget set}; her demand $x(p, m)$ is homogeneous of degree zero, and her
indirect utility is $v(p, m)$. By the envelope theorem
(\cref{thm:b1:comparative-statics:envelope}), the derivative of a \omterm{def:b1:constrained-optimisation:value-function}{value
function} with respect to a parameter is the partial derivative of the
\omterm{def:b1:constrained-optimisation:lagrangian}{Lagrangian}. A good is normal if $\partial x_i/\partial m > 0$ and inferior if
$\partial x_i/\partial m < 0$.
\end{recall}

\section{A price change split in two}\label{sec:b1:income-and-substitution-effects:split}

Take $u = x_1x_2$, income 100 and prices $(1, 1)$. The consumer buys
$(50, 50)$. The price of good 1 rises to 2, and she buys $(25, 50)$
(\cref{prop:b1:consumer-choice:cd-demand}). Her purchases of good 1 halve
for two reasons: good 1 is now dear relative to good 2, and her income of
100 buys less than before.

\begin{definition}[Substitution and income effects]\label{def:b1:income-and-substitution-effects:effects}
\index{substitution effect}\index{income effect}\index{compensation (Slutsky, Hicks)}
Let the price of good $i$ change from $p_i$ to $p_i'$. Give the consumer, at
the new prices, a \emph{compensated} income $m^c$, and let $x^c$ be her choice
then. The \emph{substitution effect} is $x^c - x(p, m)$ and the \emph{income
effect} is $x(p', m) - x^c$; they add up to the total change. With
\emph{Slutsky compensation}, $m^c$ makes the old bundle just affordable:
$m^c = m + (p_i' - p_i)x_i(p, m)$. With \emph{Hicks compensation}, $m^c$ is
the income that gives the old utility at the new prices.
\end{definition}

In the example, Slutsky compensation is $100 + (2 - 1) \times 50 = 150$, and at
prices $(2, 1)$ with income 150 she buys $x_1 = 150/4 = 37.5$. The
\omterm{def:b1:income-and-substitution-effects:effects}{substitution effect} is $37.5 - 50 = -12.5$, the \omterm{def:b1:income-and-substitution-effects:effects}{income effect}
$25 - 37.5 = -12.5$. Hicks compensation keeps utility at $50 \times 50 = 2500$:
at prices $(2, 1)$ that costs $2\sqrt{2500 \times 2} = 141.42$ (the
\omterm{def:b1:income-and-substitution-effects:expenditure}{expenditure function} below), and she buys $x_1 = 141.42/4 = 35.36$. The
\omterm{def:b1:income-and-substitution-effects:effects}{substitution effect} is then $-14.64$ and the \omterm{def:b1:income-and-substitution-effects:effects}{income effect} $-10.36$.

\begin{omfigure}
\begin{tikzpicture}
\begin{groupplot}[group style={group size=2 by 1, horizontal sep=16mm},
  omrecord, grid=none, width=0.47\linewidth, height=0.42\linewidth,
  xmin=0, xmax=80, ymin=0, ymax=110, samples=80, xlabel={$x_1$}, ylabel={$x_2$}]
\nextgroupplot[title={\small Slutsky: old bundle affordable}, width=0.47\linewidth]
  \addplot[supply] coordinates {(0,100) (80,20)};
  \addplot[supply] coordinates {(0,100) (50,0)};
  \addplot[shifted] coordinates {(20,110) (75,0)};
  \addplot[demand, domain=23:80] {2500/x};
  \addplot[demand, domain=25.6:80] {2812.5/x};
  \addplot[demand, domain=11.4:80] {1250/x};
  \addplot[only marks, mark=*, black] coordinates {(50,50) (37.5,75) (25,50)};
  \node[font=\footnotesize, anchor=north east] at (axis cs:50,50) {$A$};
  \node[font=\footnotesize, anchor=south west] at (axis cs:37.5,75) {$B$};
  \node[font=\footnotesize, anchor=north east] at (axis cs:25,50) {$C$};
\nextgroupplot[title={\small Hicks: old utility reachable}, width=0.47\linewidth]
  \addplot[supply] coordinates {(0,100) (80,20)};
  \addplot[supply] coordinates {(0,100) (50,0)};
  \addplot[shifted] coordinates {(15.71,110) (70.71,0)};
  \addplot[demand, domain=23:80] {2500/x};
  \addplot[demand, domain=11.4:80] {1250/x};
  \addplot[only marks, mark=*, black] coordinates {(50,50) (35.36,70.71) (25,50)};
  \node[font=\footnotesize, anchor=north west] at (axis cs:50,50) {$A$};
  \node[font=\footnotesize, anchor=south west] at (axis cs:35.36,70.71) {$B'$};
  \node[font=\footnotesize, anchor=north east] at (axis cs:25,50) {$C$};
\end{groupplot}
\end{tikzpicture}

{\small $u = x_1x_2$, income 100, $p_1$ from 1 to 2, $p_2 = 1$ (red lines).
Left: the dashed line has the new slope and passes through the old bundle
$A = (50, 50)$; the consumer would choose $B = (37.5, 75)$ on it, then $C = (25, 50)$
when the compensation is taken away. Right: the dashed line has the new
slope and touches the old \omterm{def:b1:preferences-and-utility:indifference-curve}{indifference curve}, at $B' = (35.36, 70.71)$.}
\end{omfigure}

The two compensations differ by a little: the Slutsky consumer can buy her
old bundle and, by moving along the new line, do better than before, so
she gets slightly more than the Hicks consumer, who is held to her old
utility. For small price changes the two effects agree to first order; the
calculus below uses the Hicks form, the data often use the Slutsky one,
since it needs only the old bundle and the prices.

\section{Compensated demand and the expenditure function}\label{sec:b1:income-and-substitution-effects:dual}

The Hicks experiment holds utility fixed and asks what it costs to reach
it. That is a problem of its own.

\begin{definition}[Expenditure function]\label{def:b1:income-and-substitution-effects:expenditure}
\index{expenditure function}
The \emph{expenditure function} is the least cost of reaching utility $u$ at
prices $p$:
\[
e(p, u) = \min_{x \ge 0}\; p_1x_1 + p_2x_2 \quad\text{subject to}\quad u(x_1, x_2) \ge u .
\]
\end{definition}

\begin{definition}[Compensated demand]\label{def:b1:income-and-substitution-effects:compensated}
\index{compensated demand}\index{Hicksian demand}
The bundle $h(p, u)$ that solves the expenditure problem is the
\emph{compensated demand}, or \emph{Hicksian demand}: what the consumer buys
at prices $p$ if her income is always adjusted to keep her at utility $u$.
\end{definition}

For $u(x) = x_1x_2$, the tangency $x_2/x_1 = p_1/p_2$ with $x_1x_2 = u$ gives
$h_1 = \sqrt{u\,p_2/p_1}$, $h_2 = \sqrt{u\,p_1/p_2}$ and
$e(p, u) = p_1h_1 + p_2h_2 = 2\sqrt{u\,p_1p_2}$. With $u = 2500$ and prices
$(2, 1)$, $e = 2\sqrt{5000} = 141.42$, the Hicks compensation of the first
section.

The two problems --- the most utility for a given income, the least
income for a given utility --- are two views of one choice. Under the
conditions of \cref{ch:b1:consumer-choice} (continuous, monotone
preferences),
\begin{equation}\label{eq:b1:income-and-substitution-effects:duality}
e\bigl(p, v(p, m)\bigr) = m,\qquad v\bigl(p, e(p, u)\bigr) = u,\qquad
h\bigl(p, v(p, m)\bigr) = x(p, m),\qquad x\bigl(p, e(p, u)\bigr) = h(p, u):
\end{equation}
the cheapest way to reach the utility that income $m$ buys costs $m$, and
the bundle is the same.

\begin{proposition}[Shephard's lemma, for the consumer]\label{prop:b1:income-and-substitution-effects:shephard}
\index{Shephard's lemma!for the consumer}
If $e$ is differentiable, $\partial e(p, u)/\partial p_i = h_i(p, u)$.
\end{proposition}

\begin{proof}
Apply the envelope theorem (\cref{thm:b1:comparative-statics:envelope},
part 2) to the minimisation, with the \omterm{def:b1:constrained-optimisation:lagrangian}{Lagrangian}
$\mathcal L = p_1x_1 + p_2x_2 - \mu\bigl(u(x) - u\bigr)$ and the parameter $p_i$:
$\partial e/\partial p_i = \partial\mathcal L/\partial p_i = x_i$, evaluated at
the solution $h(p, u)$. (A minimum is the negative of the maximum of the
negative, so the theorem applies.)
\end{proof}

The lemma says that a small rise in $p_i$ raises the cost of keeping the
consumer as well off as before by the amount she buys, $h_i\,\dd p_i$: to
first order she does not have time to substitute. On the example,
$\partial(2\sqrt{u p_1p_2})/\partial p_1 = \sqrt{u p_2/p_1} = h_1$. The producer
version, for costs, is in \cref{ch:b1:cost-minimisation}.

\begin{proposition}[Compensated demand slopes down]\label{prop:b1:income-and-substitution-effects:negative}
\index{compensated demand!slope}
$e(\cdot, u)$ is concave in $p$. Consequently, where $h$ is differentiable,
$\partial h_i/\partial p_i \le 0$.
\end{proposition}

\begin{proof}
Let $p = tp' + (1 - t)p''$ with $0 \le t \le 1$, and let $x^*$ be the cheapest
bundle reaching $u$ at $p$. Then
$e(p, u) = p\cdot x^* = t\,p'\cdot x^* + (1 - t)\,p''\cdot x^*$. Since $x^*$
reaches $u$, it costs at least $e(p', u)$ at $p'$ and at least $e(p'', u)$ at
$p''$, so $e(p, u) \ge t\,e(p', u) + (1 - t)\,e(p'', u)$: $e$ is concave in $p$.
A concave function of $p_i$ has a non-positive second derivative, and by
Shephard's lemma that derivative is $\partial h_i/\partial p_i$.
\end{proof}

\section{The Slutsky equation}\label{sec:b1:income-and-substitution-effects:slutsky}

\begin{theorem}[Slutsky equation]\label{thm:b1:income-and-substitution-effects:slutsky}
\index{Slutsky equation}
At $(p, m)$, with $u = v(p, m)$ and all functions differentiable,
\[
\frac{\partial x_i}{\partial p_j}(p, m) = \frac{\partial h_i}{\partial p_j}(p, u) - x_j(p, m)\,\frac{\partial x_i}{\partial m}(p, m).
\]
The first term is the \omterm{def:b1:income-and-substitution-effects:effects}{substitution effect}, the second the \omterm{def:b1:income-and-substitution-effects:effects}{income effect}.
\end{theorem}

\begin{proof}
By \eqref{eq:b1:income-and-substitution-effects:duality},
$h_i(p, u) = x_i\bigl(p, e(p, u)\bigr)$ for all $p$. Differentiate with respect
to $p_j$ by the chain rule:
$\partial h_i/\partial p_j = \partial x_i/\partial p_j + (\partial x_i/\partial m)(\partial e/\partial p_j)$.
By Shephard's lemma $\partial e/\partial p_j = h_j(p, u) = x_j(p, m)$. Rearrange.
\end{proof}

For the own price ($j = i$), the substitution term is negative
(\cref{prop:b1:income-and-substitution-effects:negative}). If good $i$ is
normal, the income term $-x_i\,\partial x_i/\partial m$ is negative too, and
demand slopes down: the law of demand holds for every \omterm{def:b1:consumer-choice:normal-inferior}{normal good}. The
income term is large only if the good takes a large part of the budget,
since it is weighted by $x_i$. In elasticities, with budget shares
$s_j = p_jx_j/m$ and income elasticity $\eta_i$,
\[
\varepsilon_{ij} = \varepsilon^{c}_{ij} - s_j\,\eta_i ,
\]
where $\varepsilon_{ij}$ is the elasticity of $x_i$ with respect to $p_j$ and
$\varepsilon^{c}_{ij}$ that of $h_i$ (\cref{exo:b1:income-and-substitution-effects:elasticity}).
For fuel with a budget share of $6\%$ and an income elasticity of $0.8$,
compensation changes the own-price elasticity by only
$0.06 \times 0.8 = 0.048$.

\begin{omfigure}
\begin{tikzpicture}
\begin{axis}[omrecord, grid=none, xlabel={Quantity of good 1}, ylabel={Price $p_1$},
  xmin=0, xmax=110, ymin=0, ymax=4, width=0.55\linewidth, height=0.42\linewidth, samples=80]
  \addplot[demand, domain=12.5:100] {50/x};
  \addplot[shifted, domain=25:100] {2500/x^2};
  \addplot[only marks, mark=*, black] coordinates {(50,1)};
  \node[font=\footnotesize, anchor=south west] at (axis cs:50,1) {$(50, 1)$};
  \node[font=\footnotesize, omDef, anchor=west] at (axis cs:3,1.4) {$x_1 = 50/p_1$};
  \node[font=\footnotesize, anchor=west] at (axis cs:36,2.5) {$h_1 = 50/\sqrt{p_1}$};
\end{axis}
\end{tikzpicture}

{\small Demand for good 1 with $u = x_1x_2$, $p_2 = 1$: uncompensated, income 100
(solid), and compensated at the utility 2500 that income buys at
$p_1 = 1$ (dashed). Both pass through $(50, 1)$; the compensated curve is
steeper, because it leaves out the \omterm{def:b1:income-and-substitution-effects:effects}{income effect} of a \omterm{def:b1:consumer-choice:normal-inferior}{normal good}.}
\end{omfigure}

\begin{definition}[Net substitutes and net complements]\label{def:b1:income-and-substitution-effects:net}
\index{net substitutes}\index{net complements}
Goods $i \neq j$ are \emph{net substitutes} if $\partial h_i/\partial p_j > 0$
and \emph{net complements} if $\partial h_i/\partial p_j < 0$.
\end{definition}

The definition uses \omterm{def:b1:income-and-substitution-effects:compensated}{compensated demand} on purpose. The uncompensated cross
effect $\partial x_i/\partial p_j$ mixes in the \omterm{def:b1:income-and-substitution-effects:effects}{income effect}, and can have
the opposite sign for $(i, j)$ and $(j, i)$; the compensated one cannot,
because by Shephard's lemma
$\partial h_i/\partial p_j = \partial^2e/\partial p_j\partial p_i = \partial h_j/\partial p_i$
(Young's theorem). With two goods, the two goods are always \omterm{def:b1:income-and-substitution-effects:net}{net
substitutes}: $h(p, u)$ is homogeneous of degree zero in $p$, so
$p_1\,\partial h_1/\partial p_1 + p_2\,\partial h_1/\partial p_2 = 0$, and the
own term is negative.

A seller of a good, such as the farm household of
\cref{ch:b1:consumer-choice} with its \omterm{def:b1:consumer-choice:endowment}{endowment} $\omega$, sees the income term
turn round. Her income is $p\cdot\omega$, so a rise in $p_j$ also raises it,
and the same argument gives
$\partial x_i/\partial p_j = \partial h_i/\partial p_j + (\omega_j - x_j)\,\partial x_i/\partial m$:
a net seller of good $j$ gets richer when its price rises.

\section{Giffen goods}\label{sec:b1:income-and-substitution-effects:giffen}

\begin{definition}[Giffen good]\label{def:b1:income-and-substitution-effects:giffen}
\index{Giffen good}
A good is a \emph{Giffen good} at $(p, m)$ if $\partial x_i/\partial p_i > 0$:
the consumer buys more of it when its price rises.
\end{definition}

By the \omterm{thm:b1:income-and-substitution-effects:slutsky}{Slutsky equation}, $\partial x_i/\partial p_i > 0$ needs
$-x_i\,\partial x_i/\partial m > -\partial h_i/\partial p_i \ge 0$: the good
must be inferior, and the \omterm{def:b1:income-and-substitution-effects:effects}{income effect}, weighted by the quantity bought,
must beat the \omterm{def:b1:income-and-substitution-effects:effects}{substitution effect}. In elasticities,
$s_i\,|\eta_i| > |\varepsilon^{c}_{ii}|$: the good must take a large share of
the budget. A very poor household that lives mostly on a cheap staple,
with a little of a dearer food, fits: when the staple becomes dearer, the
household can no longer afford the dearer food, and buys more of the
staple to keep its calories.

\begin{record}[Rice in Hunan]
Jensen and Miller (2008) gave randomly chosen, extremely poor households in
two Chinese provinces vouchers that lowered the price of their staple:
rice in Hunan, wheat in Gansu. They found strong evidence of Giffen
behaviour for rice in Hunan, and weaker evidence for wheat in Gansu; how
households responded depended, non-linearly, on how poor they
were. % ledger: b1:jensen2008
The prediction of the \omterm{thm:b1:income-and-substitution-effects:slutsky}{Slutsky equation} is that Giffen behaviour is
possible only where the staple is a large share of spending and is
inferior. The design --- a price change assigned at random --- answers the
\omterm{def:b1:models-and-data:identification}{identification problem} of \cref{ch:b1:models-and-data}: the subsidy moves the
price and nothing else.
\end{record}

\section{Adding up}\label{sec:b1:income-and-substitution-effects:adding-up}

The budget constraint, which holds at every price and income, restricts
how demand can respond.

\begin{proposition}[Engel and Cournot aggregation]\label{prop:b1:income-and-substitution-effects:aggregation}
\index{Engel aggregation}\index{Cournot aggregation}
With budget shares $s_i = p_ix_i/m$, income elasticities $\eta_i$ and price
elasticities $\varepsilon_{ij}$,
\[
\sum_i s_i\,\eta_i = 1\quad\text{(Engel aggregation)},\qquad
\sum_i s_i\,\varepsilon_{ij} = -s_j\quad\text{(Cournot aggregation)}.
\]
\end{proposition}

\begin{proof}
Differentiate $\sum_i p_ix_i(p, m) = m$ with respect to $m$:
$\sum_i p_i\,\partial x_i/\partial m = 1$; write
$p_i\,\partial x_i/\partial m = s_i\eta_i$. Differentiate it with respect to $p_j$:
$x_j + \sum_i p_i\,\partial x_i/\partial p_j = 0$; multiply by $p_j/m$ and write
$p_ip_j(\partial x_i/\partial p_j)/m = s_i\varepsilon_{ij}$.
\end{proof}

\omterm{prop:b1:income-and-substitution-effects:aggregation}{Engel aggregation} says that income elasticities average to one, weighted
by budget shares: if food is a necessity, something else must be a luxury.
A complete demand system --- one equation per good, estimated on household
data --- can be tested against these restrictions, against homogeneity and
against symmetry of the compensated cross effects. Adding up holds by
construction when the data satisfy the budget. Symmetry, a consequence of
utility maximisation alone, is often rejected in estimated systems, which
is evidence either against the model or against fitting one consumer to
many.

\begin{model}{The Slutsky decomposition}
\assumes The consumer of \cref{ch:b1:consumer-choice}, with smooth
preferences, so that demand and the \omterm{def:b1:income-and-substitution-effects:expenditure}{expenditure function} are
differentiable.
\predicts A non-positive own compensated slope; symmetric compensated cross
effects; for a \omterm{def:b1:consumer-choice:normal-inferior}{normal good}, an uncompensated demand that slopes down more
than the compensated one; a price rise handed back as income still lowers
the quantity bought, by the \omterm{def:b1:income-and-substitution-effects:effects}{substitution effect}.
\tested Consistent with it, though not a test of its listed predictions:
after Sweden's carbon tax and the VAT on transport fuel, emissions from
transport fell by almost $11\%$ against a comparison built from other
countries, the largest share due to the carbon tax, which came with lower
income taxes. % ledger: b1:andersson2019
Giffen behaviour, which the model allows only for an inferior staple, has
been found where it predicts: for a staple among very poor households
(Jensen and Miller 2008). % ledger: b1:jensen2008
Tests of symmetry and of the negative compensated slope need estimated
demand systems, which the second university year takes up.
\breaks When a price change is noticed differently according to its
source: in the same Swedish study, the carbon tax drew three times the
response that an equal change in the pump price would
have. % ledger: b1:andersson2019
When goods are lumpy (a household owns one car or none, and cannot
substitute at the margin).
\end{model}

\begin{record}[A carbon tax handed back]
Sweden introduced its carbon tax in 1991 at 250 kronor a tonne of carbon
dioxide, about 30 dollars; a year earlier, in March 1990, it had applied
VAT of $25\%$ to petrol and diesel. % ledger: b1:sweden-co2-tax
Both were part of a tax reform that lowered the top marginal rate of
income tax from $80\%$ to $50\%$. % ledger: b1:sweden-1991-reform
Andersson (2019) built a synthetic Sweden: a weighted average of comparable
countries that tracked Swedish transport emissions before 1990. After the
reform, the real Sweden's emissions from transport were almost $11\%$ below
the synthetic one, and the response to the carbon tax was three times
what the ordinary price elasticity of petrol demand would
predict. % ledger: b1:andersson2019
The model's prediction concerns the \omterm{def:b1:income-and-substitution-effects:effects}{substitution effect}, which the
compensation leaves in place; the size of the response is a separate
question, and the study suggests that a tax may be noticed more than an
equal change in the pump price.
\end{record}

\section*{Exercises}\label{sec:b1:income-and-substitution-effects:exercises}

\begin{exercise}[$\star$]\label{exo:b1:income-and-substitution-effects:slutsky-cd}
For $u = x_1x_2$, income 100 and prices $(1, 1)$, the price of good 1 rises to
2. Compute the Slutsky compensation, the compensated choice of good 1, and
the substitution and \omterm{def:b1:income-and-substitution-effects:effects}{income effects}.
\end{exercise}

\begin{exercise}[$\star$]\label{exo:b1:income-and-substitution-effects:inferior}
For an \omterm{def:b1:consumer-choice:normal-inferior}{inferior good}, what are the signs of the \omterm{def:b1:income-and-substitution-effects:effects}{substitution effect} and of
the \omterm{def:b1:income-and-substitution-effects:effects}{income effect} of a price rise? Can the total effect be zero?
\end{exercise}

\begin{exercise}[$\star$]\label{exo:b1:income-and-substitution-effects:read}
From the opening chart's table: Sweden emitted 2.149 tonnes per person from
road transport in 1990 and 1.554 in 2019; the mean of the ten other
countries was 1.926 and 1.928. % figdata: road_co2_points
Compute both percentage changes. Does the difference measure the effect of
the carbon tax?
\end{exercise}

\begin{exercise}[$\star$]\label{exo:b1:income-and-substitution-effects:expenditure-cd}
Derive the \omterm{def:b1:income-and-substitution-effects:compensated}{compensated demands} and the \omterm{def:b1:income-and-substitution-effects:expenditure}{expenditure function} of $u = x_1x_2$,
and check that $e\bigl(p, v(p, m)\bigr) = m$.
\end{exercise}

\begin{exercise}[$\star\star$]\label{exo:b1:income-and-substitution-effects:hicks}
On the numbers of the first exercise, compute the Hicks compensation, the
compensated choice of good 1 and the two effects. Why is the Hicks
\omterm{def:b1:income-and-substitution-effects:effects}{substitution effect} larger in size than the Slutsky one here?
\end{exercise}

\begin{exercise}[$\star\star$]\label{exo:b1:income-and-substitution-effects:shephard}
For $e(p, u) = 2\sqrt{u\,p_1p_2}$, verify Shephard's lemma for good 2 and the
symmetry $\partial h_1/\partial p_2 = \partial h_2/\partial p_1$.
\end{exercise}

\begin{exercise}[$\star\star$]\label{exo:b1:income-and-substitution-effects:fuel}
Fuel takes $6\%$ of a household's budget; its uncompensated own-price
elasticity is $-0.35$ and its income elasticity $0.8$ (invented values). The
fuel price rises by $20\%$ and the household receives a cheque that
compensates it. Predict, to first order, the fall in its fuel use.
\end{exercise}

\begin{exercise}[$\star\star$]\label{exo:b1:income-and-substitution-effects:cournot}
For Cobb--Douglas demand with share $a$ on good 1, compute $\varepsilon_{11}$,
$\varepsilon_{21}$, $\eta_1$ and $\eta_2$, and check Engel and \omterm{prop:b1:income-and-substitution-effects:aggregation}{Cournot
aggregation}.
\end{exercise}

\begin{exercise}[$\star\star$]\label{exo:b1:income-and-substitution-effects:chart}
Give three reasons why the opening chart alone cannot measure the effect
of Sweden's carbon tax on road emissions, and describe what a comparison
would need to do to measure it.
\end{exercise}

\begin{exercise}[$\star\star\star$]\label{exo:b1:income-and-substitution-effects:elasticity}
Derive the \omterm{thm:b1:income-and-substitution-effects:slutsky}{Slutsky equation} in elasticity form,
$\varepsilon_{ij} = \varepsilon^{c}_{ij} - s_j\eta_i$.
\end{exercise}

\begin{exercise}[$\star\star\star$]\label{exo:b1:income-and-substitution-effects:giffen-share}
Show that a \omterm{def:b1:income-and-substitution-effects:giffen}{Giffen good} must be inferior and must satisfy
$s_i|\eta_i| > |\varepsilon^c_{ii}|$. Use \omterm{prop:b1:income-and-substitution-effects:aggregation}{Engel aggregation} to show that the
other goods' income elasticities, weighted by their shares, then sum to more
than one. Why is a \omterm{def:b1:income-and-substitution-effects:giffen}{Giffen good} hard to find among goods with a small budget
share?
\end{exercise}

\begin{exercise}[$\star\star\star$]\label{exo:b1:income-and-substitution-effects:hunan}
In the Hunan study, households received vouchers that lowered the price of
rice. What result for rice consumption would have refuted Giffen
behaviour? Why was it essential that the vouchers were given at random?
\end{exercise}

\section{Problem: A Fuel Levy with a Cheque in the Post}\label{sec:b1:income-and-substitution-effects:problem}

\begin{problem}[{Weekend problem --- the fall in fuel use that survives the cheque}]\label{pb:b1:income-and-substitution-effects:levy}
Brenland (invented) has three kinds of household, in equal numbers, each
with preferences $u = f^{a}o^{1-a}$ over fuel $f$ (litres a year) and other
spending $o$ (dalens, price 1). Household A has income 20,000 dalens and
$a = 0.10$; B has 30,000 and $a = 0.08$; C has 50,000 and $a = 0.06$. Fuel costs
2 dalens a litre. The government adds a levy of 0.5 dalens a litre.
Quantities are given to one decimal.

\medskip\noindent\textbf{Part I --- One household, two decompositions.}
\begin{enumerate}
  \item Compute A's fuel use and other spending before the levy.
  \item Compute A's fuel use after the levy, with no compensation.
  \item Compute A's Slutsky compensated income and fuel use.
  \item Give the Slutsky substitution and \omterm{def:b1:income-and-substitution-effects:effects}{income effects}, and check that they add up.
  \item A's Hicks compensated income is $20000 \times 1.25^{0.1}$. Explain why, then compute it with A's fuel use and the two Hicks effects.
  \item Why are the Slutsky and Hicks effects so close here?
\end{enumerate}

\medskip\noindent\textbf{Part II --- Exact compensation.}
\begin{enumerate}[resume]
  \item Show that $e(p, u) = u\,p_f^{a}/\bigl(a^a(1 - a)^{1-a}\bigr)$ for these preferences, with $p_o = 1$, and that $e\bigl(p, v(p, m)\bigr) = m$.
  \item Deduce that the exact (Hicks) compensation for a levy that raises $p_f$ by the factor $1.25$ is $m(1.25^{a} - 1)$, and compute it for A, B and C.
  \item The first-order compensation is the levy times the fuel bought before it. Compute it for A, B and C and explain why it exceeds the exact one.
  \item Verify Shephard's lemma on A: $\partial e/\partial p_f = a\,e/p_f$ equals A's compensated fuel use.
  \item Compute the levy revenue from one household of each kind if no cheque is paid, and compare it with the sum of their exact compensations.
  \item Interpret the difference.
\end{enumerate}

\medskip\noindent\textbf{Part III --- The cheque.} The revenue is returned as
the same cheque $c$ to every household.
\begin{enumerate}[resume]
  \item Without the cheque, compute total fuel use of one household of each kind before and after the levy, and the fall in per cent.
  \item With the cheque, write each household's fuel demand and the condition that the cheques add up to the revenue: $3c = 0.5F$, where $F$ is total fuel use.
  \item Solve for $F$ and $c$.
  \item Compute the fall in fuel use with the cheque.
  \item If instead each household were given its exact compensation, what would total fuel use be? Compare.
  \item Explain, with the \omterm{thm:b1:income-and-substitution-effects:slutsky}{Slutsky equation}, why the cheque changes the fall so little.
\end{enumerate}

\medskip\noindent\textbf{Part IV --- Who gains.}
\begin{enumerate}[resume]
  \item Compare the cheque with each household's exact compensation. Who is better off than before the levy, and by how many dalens?
  \item Which household uses the least fuel, and which has the highest fuel share? Which of the two characteristics decides who gains?
  \item Sum the gains and losses. What does the total measure?
  \item Brenland's fuel users are told: ``the cheque gives you your money back.'' For which households is that true?
  \item Would a cheque proportional to income have changed the fall in fuel use? Use the \omterm{def:b1:consumer-choice:demand-function}{demand functions}.
  \item State the result: by how much fuel use falls with the cheque.
\end{enumerate}
\end{problem}
